\chapter{Collective Administrative Review}
\label{chap:Collective_Administrative_Review}

\section{Purpose}

Collective Administrative Review (CAR) defines how communities, institutions,
and councils exercise oversight over a CME's behaviour at scale.

Where:
\begin{itemize}
    \item ZED guards identity,
    \item P-class guards principles,
    \item Mneme guards continuity,
\end{itemize}
CAR guards the \emph{social contract}: how the CME relates to human
communities, infrastructures, and institutional commitments.

CAR is a \emph{process layer} over:
\begin{itemize}
    \item Governance Policies (Chapter~\ref{chap:Governance_Policies}),
    \item SIGIL Pathways (Chapter~\ref{chap:SIGIL_Pathways}),
    \item Ethical Alignment Council (Chapter~\ref{chap:Ethical_Alignment_Council}),
    \item Operator Interface (Chapter~\ref{chap:Review_Instructions}).
\end{itemize}

\section{Scope}

CAR applies to:
\begin{itemize}
    \item policy-level changes (P-class updates),
    \item high-impact Engram canonicalization (C-class),
    \item changes in deployment context (Where/Who/What at institutional scale),
    \item access to critical infrastructures,
    \item new or revised SIGIL pathways.
\end{itemize}

CAR does \emph{not} handle:
\begin{itemize}
    \item routine L-class reasoning,
    \item ephemeral session-level adjustments,
    \item purely local, low-risk O-class branching.
\end{itemize}

\section{Roles and Actors}

CAR operates with a set of human and institutional roles.

\subsection{Administrative Roles}

\begin{description}
    \item[Owner / Steward] Entity responsible for the CME’s existence
          (organisation, project, or community).
    \item[HITL Operator] Individual(s) performing day-to-day review in the
          Operator Interface.
    \item[Council Member] Participant in the Collective Administrative Review
          body, with voting or advisory status.
    \item[Alignment Council Liaison] Bridge role between CAR and the Ethical
          Alignment Council.
\end{description}

\subsection{CME Role}

The CME plays a bounded role in CAR:
\begin{itemize}
    \item provides Engram proposals and impact assessments,
    \item surfaces telemetry and risk analysis,
    \item does \emph{not} make final policy decisions,
    \item logs all CAR interactions in the HITL audit log.
\end{itemize}

\section{Review Classes}

CAR distinguishes review by impact level.

\subsection{Class~A: Kernel / Charter Level}

Affects:
\begin{itemize}
    \item ZED Core (identity kernel),
    \item P-class charters,
    \item foundational alignment principles.
\end{itemize}

Requirements:
\begin{itemize}
    \item quorum of Council Members,
    \item explicit involvement of Ethical Alignment Council,
    \item SIGIL-level authentication for final approval,
    \item full archival in C-class and audit log.
\end{itemize}

\subsection{Class~B: Policy / Governance Level}

Affects:
\begin{itemize}
    \item Governance Policies,
    \item access scopes and deployment contexts,
    \item new classes of permissible or forbidden Engrams,
    \item GEL admission criteria.
\end{itemize}

Requirements:
\begin{itemize}
    \item multi-person review (at least two approvers),
    \item Council notification and review window,
    \item audit log entry with rationale.
\end{itemize}

\subsection{Class~C: Operational Level}

Affects:
\begin{itemize}
    \item specific Engram canonicalization decisions,
    \item tag-level changes that touch P- or C-class,
    \item creation or modification of high-impact SIGIL pathways.
\end{itemize}

Requirements:
\begin{itemize}
    \item single-operator approval with optional escalation,
    \item standard HITL logging,
    \item rollback plan in case of later reversal.
\end{itemize}

\section{Input Channels}

CAR receives input from:

\begin{itemize}
    \item \textbf{Operator Approval Queue} (Chapter~\ref{chap:Operator_Approval_Queue}),
    \item automated risk detectors in the CME harness,
    \item human stakeholders submitting policy questions,
    \item periodic reviews triggered by time or usage thresholds.
\end{itemize}

Each input is normalized into a \emph{Review Dossier}:
\begin{itemize}
    \item W5H summary of the proposed change,
    \item impact assessment,
    \item class (A/B/C) assignment,
    \item recommended decision paths.
\end{itemize}

\section{Review Workflow}

\subsection{Step 1: Intake and Classification}

\begin{enumerate}
    \item Dossier enters CAR queue.
    \item Automated pre-checks determine:
    \begin{itemize}
        \item affected classes (C/L/M/S/O/P),
        \item linkage to P-class charters,
        \item potential impact on $Inv(E)$ or ZED.
    \end{itemize}
    \item Assign review class (A/B/C).
\end{enumerate}

\subsection{Step 2: Deliberation}

\begin{itemize}
    \item For Class~A:
    \begin{itemize}
        \item joint session with Ethical Alignment Council,
        \item consideration of long-term identity impact,
        \item may require revisions to other governance artefacts.
    \end{itemize}
    \item For Class~B:
    \begin{itemize}
        \item Council review (asynchronously or synchronously),
        \item comments and requested modifications recorded,
        \item potential escalation to Class~A if scope expands.
    \end{itemize}
    \item For Class~C:
    \begin{itemize}
        \item operator review against existing policies,
        \item optional consult with Council Liaison,
        \item decision within defined SLA window.
    \end{itemize}
\end{itemize}

\subsection{Step 3: Decision and Encoding}

Possible outcomes:
\begin{itemize}
    \item \textbf{Approve}: change enacted; Engram or policy updated.
    \item \textbf{Approve with modifications}: applicant revises and resubmits.
    \item \textbf{Reject}: change denied with rationale logged.
    \item \textbf{Escalate}: moved to higher class of review.
\end{itemize}

Encoding:
\begin{itemize}
    \item Updates to Governance Policies or P-class charters,
    \item changes to allowed/forbidden Engram types,
    \item updates to SIGIL pathways and operator instructions.
\end{itemize}

\section{Interaction with Other HITL Components}

\subsection{Ethical Alignment Council}

Class~A reviews always invoke the Ethical Alignment Council.  
Council decisions:
\begin{itemize}
    \item can veto or demand changes to proposals,
    \item may declare ``hard constraints'' that CAR must respect.
\end{itemize}

\subsection{SIGIL Pathways}

High-impact CAR decisions may:
\begin{itemize}
    \item create new SIGIL pathways for sensitive operations,
    \item revoke or modify existing SIGIL keys,
    \item change the mapping between symbolic rituals and technical actions.
\end{itemize}

\subsection{Operator Interface and Audit Log}

All CAR decisions must be:
\begin{itemize}
    \item visible in the Operator Interface,
    \item recorded in the HITL audit log (Chapter~\ref{chap:HITL_Audit_Log}),
    \item cross-referenced with affected Engrams and policies.
\end{itemize}

\section{Safeguards and Failure Modes}

CAR includes safeguards against:
\begin{itemize}
    \item unilateral policy changes by a single operator,
    \item silent drift in governance artefacts,
    \item inconsistent decisions across time or contexts.
\end{itemize}

Failure modes:
\begin{description}
    \item[Deadlock] No decision reached in reasonable time.  
          Mitigation: escalation to a higher council or external authority.
    \item[Inconsistency] Decisions conflict with prior rulings or charters.  
          Mitigation: forced reconciliation, supersession records, public log.
    \item[Bypass] Actions taken outside CAR processes.  
          Mitigation: audit log anomaly detection, SIGIL revocation, emergency halt.
    \item[Misalignment] CAR outcome conflicts with P-class or ZED.  
          Mitigation: automatic escalation to Ethical Alignment Council.
\end{description}

\section{Summary}

Collective Administrative Review provides:
\begin{itemize}
    \item a structured, multi-layered process for governing the CME,
    \item clear separation between operational decisions and policy-level changes,
    \item institutional memory for why changes were made,
    \item and a social contract layer over technical and cognitive safeguards.
\end{itemize}

It ensures that the CME’s evolution is not just technically safe, but also
contextually legitimate for the communities that live with it.
